\section{Introduction}

Deep neural networks now reach high accuracy on public chest X-ray (CXR) classification benchmarks, and architectures such as DenseNet121 have been reported to perform comparably to radiologists on specific detection tasks~\cite{rajpurkar2017chexnet}. However, high test accuracy does not show \emph{where} a model finds its evidence. Models may use text annotations, image borders, patient-positioning cues or scanner-specific intensity patterns that happen to correlate with the label in the training data. Such a model can therefore look reliable on a benchmark and still fail on patients whose images carry the same non-clinical cues. Suppose, for example, that most disease-positive training images come from one hospital that marks its radiographs with a particular text annotation and border style. A radiograph of a healthy patient from that same hospital can then be classified as diseased, because the model recognises the hospital's markings rather than any pathology in the lungs. The reverse can also happen: a diseased patient imaged at a different site may be missed. In a clinical setting, both kinds of error are harmful.

Several studies have documented this failure mode. Zech et al.~\cite{zech2018variable} showed that a pneumonia classifier could identify the hospital system an image came from, and that its performance dropped on radiographs from external sites. DeGrave et al.~\cite{degrave2021shortcuts} found that COVID-19 detectors trained on multi-source datasets responded to laterality markers, text tokens and image borders rather than only to pulmonary findings. Oakden-Rayner et al.~\cite{oakdenrayner2020hidden} showed that a single aggregate metric can hide clinically important failures in subsets of the data. Geirhos et al.~\cite{geirhos2020shortcut} describe this general behaviour as \emph{shortcut learning}: a model takes a decision rule that works on the benchmark but does not reflect the intended task. Shortcuts are hard to detect because two models with nearly identical accuracy may base their predictions on very different image regions.

The COVID-19 Radiography Database~\cite{chowdhury2020covidradiography,rahman2021enhancement}, used in this work, is a typical example of this risk. Its four classes (COVID-19, Normal, Lung Opacity and Viral Pneumonia) were collected from different public sources, so acquisition properties may be confounded with the diagnostic label. We do not assume that such shortcuts have already been confirmed in this dataset. Instead, we ask a narrower and testable question: \emph{can lung-mask supervision, used only during training, move a classifier's decision evidence into the lung fields without reducing accuracy, and does the same design transfer across architectures?} The dataset provides a lung mask for every image, which makes it possible to supervise spatial attention directly. It also lets us measure the result with Grad-CAM~\cite{selvaraju2017gradcam} Energy-Inside-Lung (EIL), the fraction of Grad-CAM energy that falls inside the lung mask.

Earlier approaches to this problem have limitations. Attention modules such as CBAM~\cite{woo2018cbam} learn spatial weighting only from the classification loss, so nothing directs them to anatomy. Segmentation-first pipelines~\cite{bassi2022segmentation} and mask-guided pooling frameworks~\cite{mamun2026beyond} use the lungs explicitly, but they need a segmentation stage at inference time or change the backbone substantially. Debiasing objectives~\cite{luo2022pseudobias} avoid masks but do not specify which region the model should use. Most prior evaluations also study a single backbone, so it is unclear whether a proposed fix is a property of the method or of one particular architecture.

In this paper we propose a lightweight \emph{Lung-Region Attention} module. It is inserted between a pretrained backbone and its classification head, and it is supervised with soft lung-mask targets through an auxiliary loss. The mask is needed only to compute that loss during training, so the deployed classifier takes the CXR image alone. We develop the module on DenseNet121~\cite{huang2017densenet} and then transfer it, under the same data split, seed policy and evaluation protocol, to ResNet50~\cite{he2016resnet}, EfficientNet-B0~\cite{tan2019efficientnet} and ViT-Base/16~\cite{dosovitskiy2021vit}. On each backbone we run a controlled six-arm ablation that separates mask guidance from attention gating and compares a residual gate, $f' = f \odot (1 + a)$, with a purely multiplicative gate, $f' = f \odot a$. The arms also include an optional background-suppression term and a CBAM spatial-attention comparison. For each backbone, we select the winning arm with a pre-registered rule: an arm qualifies only if McNemar's test detects no accuracy degradation and a paired Wilcoxon test confirms an EIL gain. Winners are then re-evaluated over three matched seeds.

The three convolutional backbones show a consistent result: mask-guided gating moves Grad-CAM evidence into the lungs without a meaningful loss of accuracy, but the best form of gate depends on the backbone. On DenseNet121, the residual gate raises EIL by $7.05 \pm 0.60$ percentage points (pp) across three seeds, with no statistically significant change in accuracy. On ResNet50 and EfficientNet-B0, the multiplicative gate wins. It raises EIL by about $29.9$\,pp on ResNet50 and by $11.3 \pm 2.6$\,pp on EfficientNet-B0, and on EfficientNet-B0 it also improves accuracy by $0.76 \pm 0.08$\,pp. The module adds at most 2.6\% parameters and 1.3\% FLOPs to any of the convolutional backbones. We also report a negative result. On EfficientNet-B0, a counterfactual test that zeroes, shuffles or noises the non-lung pixels changes the predicted class for a large majority of test images in \emph{every} arm, including the attention-gated winner. Better Grad-CAM localization therefore does not, on its own, show that a model relies less on background information. EIL should be read as a measure of where evidence is placed, not of causal reliance~\cite{adebayo2018sanity}.

The main contributions of this paper are:
\begin{itemize}
\item \textbf{A mask-free-at-inference lung attention module.} The module is a lightweight attention gate supervised with soft lung-mask targets that are pooled to the feature resolution. It adds negligible inference cost and needs no segmentation mask at deployment.
\item \textbf{A cross-architecture, six-arm ablation.} We run the same six-arm ablation on DenseNet121, ResNet50, EfficientNet-B0 and ViT-Base/16 under one shared protocol. It separates the effects of mask guidance, attention gating, the choice between residual and multiplicative gating, and background suppression, and it includes a CBAM baseline.
\item \textbf{A statistically gated winner-selection protocol.} Paired McNemar and Wilcoxon tests decide each backbone's winner on a no-accuracy-cost, positive-EIL-gain criterion, and the winners are confirmed across three matched seeds.
\item \textbf{Localization and counterfactual evidence reported side by side.} We combine Grad-CAM EIL, attention--mask agreement, efficiency measurements and background-perturbation tests. This shows where attention supervision helps and where it does not yet reduce background sensitivity.
\end{itemize}

The rest of the paper is organized as follows. Section~2 reviews related work on shortcut learning and anatomically guided attention in CXR analysis. Section~3 describes the Lung-Region Attention module and its training objective. Section~4 presents the dataset, the training protocol and the evaluation metrics. Section~5 reports the per-backbone ablations, the multi-seed confirmation, and the efficiency and counterfactual analyses. Section~6 discusses limitations, and Section~7 concludes.